\documentclass[10pt,a4paper]{amsart}
\usepackage[margin=19mm]{geometry}
\usepackage{amsmath,amssymb,mathtools}
\usepackage[scr=rsfs,cal=euler]{mathalfa}
\usepackage{xfrac}
\usepackage[unicode,hidelinks,bookmarks=false]{hyperref}
\newcommand{\ie}{i.e.\ }
\newcommand{\cf}{cf.\ }
\newcommand{\resp}{respectively\ }
\newcommand{\Subsection}{\subsection}
\providecommand{\nobreakdash}{\nobreak\hbox{-}}
\setlength{\emergencystretch}{2em}
\multlinegap0pt
\usepackage{bm}
\usepackage{bbm}
\usepackage{dsfont}
\usepackage{xspace}
\usepackage{cancel}
\usepackage{mathtools}
\usepackage{amsthm}
\makeatletter
\@ifundefined{lemma}{\newtheorem{lemma}{Lemma}}{}
\makeatother
\newcommand{\lfpfacef}{\mathcal{F}}
\newcommand{\lfp}{\mathcal{P}}
\title[Perfect Copositive Matrices]{Perfect Copositive Matrices}
\author{Valentin Dannenberg, Achill Sch\"urmann}
\date{Source: arXiv:2303.17310v4 (CC BY 4.0)}
\begin{document}
\maketitle

\noindent\emph{Source and licence.} Perfect Copositive Matrices. Version arXiv:2303.17310v4 (CC BY 4.0), distributed under the Creative Commons Attribution 4.0 International licence, as recorded in the arXiv licence metadata for this exact version and shown on the abstract page https://arxiv.org/abs/2303.17310v4. Full rights notice: licenses/arxiv-2303.17310-CC-BY-4.0.txt.\par\smallskip

\noindent\textit{Editorial scope.} Counted unit: the Lemma whose source label is \texttt{lem:faceanalog} in the article (source0.tex lines 1251--1255, source label \texttt{lem:faceanalog}), delivered together with the article's own complete original proof of it (source0.tex lines 1256--1283, one \texttt{proof} environment of 28 lines). No mathematics is written, repaired or re-derived: statement and proof are the article's own text line for line. Required context retained uncounted: none. Every definition, display and label of those lines is retained unchanged. Omitted from this package and not redistributed: 1--1250, 1284--1611. Nothing inside a retained excerpt is omitted or altered: every retained body is delivered exactly as the article's lines are written, so no omission receipt of any kind accompanies this package. Citations: the retained excerpts carry no citation command at all, so no bibliography entry of the article is transcribed and no reference is redistributed. Reference adaptation: none was needed and none was made; every reference inside the delivered unit resolves inside it, so no unresolved reference or citation is produced. Printed slips kept literal: no printed word, symbol, hypothesis or number of the counted statement or of its proof was altered. Layout and class: the article's own class is \texttt{commat} with packages including amsmath, amssymb, amsfonts, color, enumitem, graphicx, hyperref, booktabs, microtype, standalone, tikz; the wrapper is the lane's pinned \texttt{amsart} editorial wrapper at 10pt on A4, which restates the theorem-like environments the retained text needs and adds the article's own macro definitions that the retained text needs (\texttt{\textbackslash lfp}, \texttt{\textbackslash lfpfacef}), with extra wrapper packages (bm, bbm, dsfont, xspace, cancel, mathtools). The delivered numbering is the unit's own. Assets: no retained span contains a figure environment, a graphics-inclusion command, a PDF-page inclusion, a table or a TikZ picture; no image file is copied into this package, the package declares no asset dependency and no image file is redistributed with it. Licence: version arXiv:2303.17310v4 of this article is distributed under the Creative Commons Attribution 4.0 International licence, as recorded in the arXiv licence metadata for this exact version and shown on the abstract page https://arxiv.org/abs/2303.17310v4. A full rights notice is delivered at licenses/arxiv-2303.17310-CC-BY-4.0.txt. Characters: every retained excerpt is pure ASCII, so the delivered engine, XeLaTeX with its default Latin Modern text font, reports no missing character.
\begin{lemma}\label{lem:faceanalog}
	Let \(\lfp=\bigcap_{i=1}^{\infty}H^{\leq}(a_i,\beta_i) \) be a~locally finite polyhedron.
	Then \(\lfpfacef \) is a~face of \(\lfp \) if and only if there is a~finite index set $I$,
	such that \(\lfpfacef=\lfp\cap\bigcap_{i\in I}H(a_i,\beta_i) \).
\end{lemma}

\begin{proof}
	Note that \(H=\left\{ x:\ \langle\sum_{i\in I}a_i,x\rangle = \sum_{i\in I}\beta_i \right\} \)
	is a~supporting hyperplane of \(\lfp \).
    As \(\lfpfacef=\lfp\cap H \), we see that \(\lfpfacef \) is a~face of \(\lfp \).

	On the other hand, let \(\lfpfacef \) be a~face of \(\lfp \).
	Then there is some hyperplane $H$ with \(\lfpfacef = \lfp\cap H \) and
    \(\lfp\subseteq H^{\leq} \).
	By choosing $x$ in the relative interior of \(\lfpfacef \), we find an affine subspace $A$
	with \(A \subseteq H \) and \(\lfpfacef = \lfp\cap A \).
	Now take some arbitrary polytope $C$ containing $x$ in its interior.
	Then, the intersection \(C\cap\lfp \) is a~polytope by assumption, and further
	
\[
		C\cap\lfpfacef=C\cap\lfp\cap A \text{ is a~face of } C\cap\lfp.
\]
	As it is a~polytope it can be described as an intersection of finitely many halfspaces, some of
	them among the \(H^{\leq}(a_i, \beta_i) \).
	However the halfspaces with bounding hyperplanes containing $A$ can be chosen solely from the
	\(H^{\leq}(a_i, \beta_i) \).
	In other words
	
\[
		A = \bigcap_{i\in I}H(a_i, \beta_i) \text{ with a~suitable finite index set } I.
\]
	Therefore
	\(\lfpfacef = \lfp\cap\bigcap_{i\in I}H(a_i, \beta_i). \)
\end{proof}

\end{document}
